% =========================
% sections/introduction.tex
% =========================
\section{Introduction}

Healthcare systems worldwide are undergoing a fundamental digital transformation. Adoption of interoperable electronic health record (EHR) systems, health information exchanges (HIEs), and broader eHealth infrastructure has become a central priority for governments, healthcare organizations, and technology vendors. In the United States, the Health Information Technology for Economic and Clinical Health (HITECH) Act of 2009 catalyzed national EHR-adoption incentives exceeding \$35 billion~\cite{adlermilstein2015hitech35b,onc2025reports}, while national digital health strategies in the United Kingdom, Estonia, Australia, Canada, and across the European Union have driven similar large-scale investment in interoperable health information infrastructure.

Despite this sustained attention and investment, adoption outcomes remain strikingly uneven. Some systems—such as Estonia's national eHealth platform—have achieved near-comprehensive coverage, high clinician engagement, and measurable care-coordination gains~\cite{estonia2026ehealth,parv2012estoniahie}. Others, including heavily funded U.S. programs such as the Department of Veterans Affairs (VA) and Department of Defense (DoD) joint EHR modernization, continue to face usability challenges, implementation delays, and user dissatisfaction years after deployment~\cite{gao2025vaehr,gao2024dodehr}. This paradox—substantial investment yielding inconsistent outcomes—motivates the central inquiry of this study.

Adoption science offers several relevant frameworks. The DeLone and McLean IS Success Model~\cite{delone2003model} identifies system quality, information quality, use, and net benefits as critical success dimensions; the Technology-Organization-Environment (TOE) framework~\cite{tornatzky1990processes} situates adoption within technological readiness, organizational capacity, and environmental pressures. Individually illuminating, neither was developed for the specific complexities of healthcare—a domain of high regulatory density, sensitive data governance, deeply embedded clinical workflows, and multi-stakeholder governance. Prior systematic reviews have tended to examine isolated determinants—usability, interoperability standards, policy mandates, or organizational leadership—without integrating them into a unified, operationalizable model applicable across diverse contexts~\cite{eden2016barriers,kruse2016adoption,aguirre2019ehr}.

This study addresses that gap and extends a sustained line of inquiry. It builds on the author's prior qualitative research~\cite{adams2020strategies}, which identified adoption strategies from interviews with senior healthcare IT leaders, by synthesizing the broader systematic-review literature and validating the resulting framework against four publicly documented national implementations. I develop and validate the \textbf{U-T-I-O Model}—a structured framework integrating \textbf{U}ser Readiness, \textbf{T}echnical Interoperability, \textbf{I}nstitutional Alignment, and \textbf{O}rganizational Execution as the four interdependent determinants of successful eHealth systems adoption. By incorporating the institutional and policy dimensions that distinguish healthcare from other IS contexts, the model gives leaders, system architects, and policymakers a practical diagnostic for assessing adoption readiness and identifying domain-specific gaps. Its scope is explicitly bounded to \textit{implementation readiness and maturity}: it characterizes the organizational, technical, institutional, and human conditions under which eHealth systems are positioned for sustainable adoption, not direct measurement of patient clinical outcomes—a relationship identified as the primary direction for future research and addressed in the Conclusion.

The study contributes in three ways: it consolidates fragmented adoption determinants into a four-domain model specific to interoperable healthcare systems; it validates the model against four publicly documented national implementation cases; and it distinguishes implementation-readiness validity from clinical-outcome validity, establishing a pathway for future U-T-I-O-OM validation. The remainder of the article reviews the theoretical foundations (Section~2), presents the research questions (Section~3), describes the methodology (Section~4), reports findings (Section~5), discusses implications (Section~6), and addresses limitations and conclusions (Sections~7--8).
